\documentclass[11pt,a4paper]{article}
\usepackage[utf8]{inputenc}
\usepackage[T1]{fontenc}
\usepackage{amsmath,amssymb,amsthm}
\usepackage{geometry}
\usepackage{hyperref}
\usepackage{booktabs}

\geometry{margin=1.2in}
\hypersetup{colorlinks=true,linkcolor=blue!70!black,citecolor=green!50!black}

\theoremstyle{plain}
\newtheorem{theorem}{Theorem}
\newtheorem{proposition}[theorem]{Proposition}
\theoremstyle{definition}
\newtheorem{definition}[theorem]{Definition}
\newtheorem{remark}[theorem]{Remark}

\newcommand{\BISH}{\textsf{BISH}}
\newcommand{\WKL}{\textsf{WKL}_0}
\newcommand{\WLPO}{\textsf{WLPO}}
\newcommand{\CLASS}{\textsf{CLASS}}
\newcommand{\Q}{\mathbb{Q}}
\newcommand{\Z}{\mathbb{Z}}
\newcommand{\Sha}{\mathrm{III}}
\newcommand{\rank}{\mathrm{rank}}

\title{\textbf{Paper 97: BSD Landscape Survey --- Null Finding}\\[6pt]
\large Internal Note\\[4pt]
\normalsize\textit{CRM Programme, Papers 50--98}}
\author{Paul Chun-Kit Lee}
\date{March 2026}

\begin{document}
\maketitle

\begin{abstract}
We report a null finding from a systematic computational and analytical survey of the BSD conjecture landscape, seeking CRM bifurcation analogous to the Hodge detection result (Papers~86--90). Three directions were explored: explicit 2-descent excision of Sha finiteness, $p$-adic Gross--Zagier bypass of Archimedean $L$-values, and parameter-space bifurcation via root number stratification. All three fail. The null finding is itself informative: the moduli dimension argument (elliptic curves on the 1-dimensional $j$-line vs.\ genus-4 Jacobians in 9-dimensional moduli) explains structurally why BSD cannot bifurcate. One positive extraction: the rank/Sha decoupling theorem ($\rank$ verification $\in \BISH$, $\Sha$ finiteness $\in \CLASS$). This note documents the computational survey (2214 curves, conductor $\leq 500$) and the three analytical dead ends for the programme record.
\end{abstract}

\section{Motivation}

Papers~86--90 established that the Hodge conjecture for Weil fourfolds exhibits CRM bifurcation: detection is $\WLPO$ generically but $\BISH$ on the palindromic locus, where the Kani--Rosen splitting provides an algebraic shortcut. The natural question: does BSD exhibit analogous bifurcation? Specifically, can $\Sha$ finiteness be proved constructively for a nontrivial sublocus of elliptic curves?

\section{Computational Survey}

A Python survey (SageMath/LMFDB interface) examined 2214 elliptic curves with conductor $N \leq 500$:

\begin{center}
\begin{tabular}{@{}lrl@{}}
\toprule
\textbf{Invariant} & \textbf{Count/Range} & \textbf{Observation} \\
\midrule
Total curves & 2214 & Cremona database \\
Root number $w = +1$ & $\sim 1107$ & Analytic rank even \\
Root number $w = -1$ & $\sim 1107$ & Analytic rank odd \\
Analytic rank 0 & majority ($w = +1$) & $L(E,1) \neq 0$, Sha finite (Kolyvagin) \\
Analytic rank 1 & majority ($w = -1$) & $L'(E,1) \neq 0$, Sha finite (Gross--Zagier--Kolyvagin) \\
Analytic rank $\geq 2$ & rare & BSD unproved \\
Anomalous BSD behavior & \textbf{0} & No curves found \\
\bottomrule
\end{tabular}
\end{center}

No anomalies detected. No curves exhibited behavior suggestive of constructive shortcuts. The root number partition $w = \pm 1$ is the only structural stratification, and it is discrete (not continuous).

\section{Three Analytical Directions: All Fail}

\subsection{Direction 1: Explicit 2-Descent Excision}

\textbf{Idea.} 2-descent computes $\rank E(\Q)$ (via the 2-Selmer group, which is finite Galois cohomology: $\BISH$). If $\Sha[2] = 0$ could be leveraged to prove $|\Sha| < \infty$, this would give a $\BISH$ path to BSD.

\textbf{Failure.} $\Sha[2] = 0$ says nothing about odd-primary components $\Sha[p^\infty]$ for $p \geq 3$. Proving $|\Sha| < \infty$ requires annihilating infinitely many $p$-primary components simultaneously. Kolyvagin's Euler system is the only known method, and it requires:
\begin{enumerate}
\item[(a)] A Heegner point of infinite order (from Gross--Zagier: Archimedean $L'(E,1)$).
\item[(b)] Chebotarev density theorem (analytic continuation of Artin $L$-functions).
\end{enumerate}
Both are irreducibly $\CLASS$.

\textbf{Verdict: FATAL GAP.} Rank/$\Sha$ decoupling is sharp. No bridge from 2-descent to full Sha finiteness.

\subsection{Direction 2: $p$-adic Gross--Zagier}

\textbf{Idea.} Coleman integration is $\BISH$ (algebraic de Rham, no Archimedean place). The $p$-adic Gross--Zagier formula relates $p$-adic $L'$-values to Heegner points via $p$-adic integration. Could this bypass the Archimedean $L$-value?

\textbf{Failure.} The comparison theorem between $p$-adic and Archimedean $L$-functions reintroduces the complex period $\Omega_+(E)$. To conclude anything about the \emph{actual} $L$-function (which BSD concerns), one must compare the $p$-adic $L$-function with the complex one. This comparison passes through the Betti--de Rham comparison isomorphism: integration of algebraic differential forms over topological cycles on $E(\mathbb{C})$. This is exactly the Archimedean obstruction identified in Paper~98.

\textbf{Verdict: CIRCULAR.} The $p$-adic detour avoids $\mathbb{R}$ locally but reintroduces it at the comparison step. False accounting of logical cost.

\subsection{Direction 3: Parameter-Space Bifurcation}

\textbf{Idea.} Could special families of elliptic curves (CM curves, Heegner-friendly conductors, etc.) admit constructive Sha finiteness proofs, analogous to how the palindromic locus admits constructive Hodge detection?

\textbf{Failure: Moduli dimension argument.} Genus-4 Jacobians live in a 9-dimensional moduli space ($\mathcal{M}_4 \hookrightarrow \mathcal{A}_4$). The palindromic sublocus $\{a = c\}$ is a continuous family of positive codimension, providing geometric degrees of freedom for exceptional behavior.

Elliptic curves are isolated points on the 1-dimensional $j$-line. There are no continuous families with tunable parameters that could create an exceptional sublocus. CM curves are isolated points (finitely many $j$-invariants for each discriminant). The root number $w = \pm 1$ is a discrete invariant, not a continuous parameter.

\textbf{Verdict: DEFINITIVE NEGATIVE.} Moduli dimension 1 precludes bifurcation. The Hodge result required moduli dimension $\geq 3$ (palindromic is codimension 3 in $\mathcal{A}_4$). BSD cannot replicate this.

\section{Positive Extraction: Rank/Sha Decoupling}

\begin{theorem}[Rank/$\Sha$ Decoupling]
The two components of BSD have different CRM signatures:
\begin{enumerate}
\item[(i)] \textbf{Rank verification} ($\BISH$): Given candidate points $P_1, \ldots, P_r \in E(\Q)$, verifying linear independence and computing $\rank\langle P_1, \ldots, P_r \rangle$ is $\BISH$. The 2-Selmer group is finite Galois cohomology with $\Z/2\Z$ coefficients.
\item[(ii)] \textbf{$\Sha$ finiteness} ($\CLASS$): Proving $|\Sha(E/\Q)| < \infty$ requires Kolyvagin's Euler system (Heegner points from Gross--Zagier, annihilation via Chebotarev), which is irreducibly $\CLASS$.
\end{enumerate}
\end{theorem}

This theorem is incorporated into Paper~98 (Calibration Theorem III, \S8).

\section{Structural Explanation via Paper 98}

The null finding is fully explained by the Archimedean Obstruction Thesis of Paper~98. The BSD conjecture concerns the special value $L(E, 1)$ (or $L'(E, 1)$), which involves:

\begin{enumerate}
\item[(a)] The Archimedean $\Gamma$-factor: $\CLASS$.
\item[(b)] The period integral $\Omega_+(E) = \int_{E(\mathbb{R})} \omega$: $\CLASS$ (integration over $\mathbb{R}$).
\item[(c)] The Mellin transform defining $L(E, s)$: $\CLASS$ (integral over $\mathbb{R}_{>0}$).
\end{enumerate}

Every path to BSD passes through the Archimedean place. Unlike the Hodge conjecture, where the detection problem (``is this class in $H^{2,2}$?'') can sometimes be answered by algebraic methods on special fibers, the BSD conjecture asks about a specific complex number ($L(E,1)$ or its derivative), and there is no algebraic substitute for this analytic quantity.

\section{Conclusion}

The BSD landscape survey yields a clean null result: no CRM bifurcation exists for BSD, no constructive sublocus for $\Sha$ finiteness, no parameter-space structure enabling algebraic shortcuts. The null finding is explained by the moduli dimension argument (dimension 1 precludes bifurcation) and confirmed by the Archimedean Obstruction Thesis (every BSD proof path passes through $\mathbb{R}$).

The extractable content --- the rank/$\Sha$ decoupling theorem --- has been incorporated into Paper~98 as one of three calibration theorems.

\bigskip
\noindent\textit{This is an internal programme note documenting a null finding. It is not intended for independent publication but serves as a record for the Papers 50--98 CRM programme.}

\end{document}
